% Title: Tensor Product Representations on Riemannian Structures
% Author: Amlal El Mahrouss (A. E. Mahrouss)

\documentclass[10pt]{article}

%% ---- Core Packages ----------------------------------------
\usepackage[T1]{fontenc}
\usepackage[utf8]{inputenc}
\usepackage{lmodern}
\usepackage{microtype}

%% ---- Page Geometry ----------------------------------------
\usepackage[
  top=1in, bottom=1in,
  left=0.75in, right=0.75in,
  columnsep=0.25in
]{geometry}

%% ---- Mathematics ------------------------------------------
\usepackage{amsmath, amssymb, amsthm, amsfonts}
\usepackage{mathtools}
\usepackage{bm}
\usepackage{textcomp}
\usepackage{esint}

%% ---- Figures & Tables -------------------------------------
\usepackage{graphicx}
\usepackage{booktabs}
\usepackage{multirow}
\usepackage{array}
\usepackage{caption}
\usepackage{subcaption}
\usepackage{float}

%% ---- Code Listings ----------------------------------------
\usepackage{listings}
\usepackage{xcolor}
\usepackage{algorithm}
\usepackage{algpseudocode}

\lstdefinestyle{codestyle}{
  backgroundcolor=\color{gray!8},
  basicstyle=\ttfamily\footnotesize,
  breakatwhitespace=false,
  breaklines=true,
  captionpos=b,
  commentstyle=\color{green!50!black},
  keywordstyle=\color{blue}\bfseries,
  stringstyle=\color{orange!70!black},
  numberstyle=\tiny\color{gray},
  numbers=left,
  numbersep=5pt,
  showstringspaces=false,
  frame=single,
  rulecolor=\color{gray!40},
  tabsize=2
}
\lstset{style=codestyle}

%% ---- Hyperlinks -------------------------------------------
\usepackage[
  colorlinks=true,
  linkcolor=blue!70!black,
  citecolor=green!50!black,
  urlcolor=blue!60!black
]{hyperref}
\usepackage{url}

%% ---- Bibliography -----------------------------------------
\usepackage[numbers, sort&compress]{natbib}

%% ---- Miscellaneous ----------------------------------------
\usepackage{enumitem}
\usepackage{siunitx}
\usepackage{cleveref}

%% ---- Theorem Environments ---------------------------------
\theoremstyle{plain}
\newtheorem{theorem}{Theorem}[section]
\newtheorem{lemma}[theorem]{Lemma}
\newtheorem{corollary}[theorem]{Corollary}
\newtheorem{proposition}[theorem]{Proposition}

\theoremstyle{definition}
\newtheorem{definition}[theorem]{Definition}
\newtheorem{example}[theorem]{Example}

\theoremstyle{remark}
\newtheorem{remark}[theorem]{Remark}

%% ---- Custom Commands --------------------------------------
\newcommand{\R}{\mathbb{R}}
\newcommand{\N}{\mathbb{N}}
\newcommand{\Z}{\mathbb{Z}}
\newcommand{\E}{\mathbb{E}}
\newcommand{\Prob}{\mathbb{P}}
\newcommand{\norm}[1]{\left\lVert #1 \right\rVert}
\newcommand{\abs}[1]{\left\lvert #1 \right\rvert}
\newcommand{\inner}[2]{\left\langle #1,\, #2 \right\rangle}
\newcommand{\ie}{\textit{i.e.}\xspace}
\newcommand{\eg}{\textit{e.g.}\xspace}
\newcommand{\etal}{\textit{et al.}\xspace}

%% ---- Caption Formatting -----------------------------------
\captionsetup{
  font=small,
  labelfont=bf,
  format=hang,
  justification=justified
}

\usepackage{lettrine}

%% ---- Header / Footer --------------------------------------
\usepackage{fancyhdr}
\pagestyle{fancy}
\fancyhf{}
\fancyhead[R]{\small\thepage}
\renewcommand{\headrulewidth}{0.4pt}

%% ============================================================
%%  FRONT MATTER
%% ============================================================

\title{%
	\vspace{-1.5em}%
	\large\textbf{Tensor Product Representations on Riemannian Structures.}%
}
\author{By Amlal El Mahrouss}
\date{\today}

%% ============================================================
\begin{document}
%% ============================================================

\maketitle
\tableofcontents
\newpage

%% ============================================================
\section{Abstract:}
%% ============================================================

\lettrine[lines=1, lhang=0.1, loversize=0.2]{T}HE research paper covers and explores several mathematical concepts of Riemannian Manifolds and Structures in Abstract Algebra and Topology. We prove that:\\
\begin{equation}
  \int_{\mathcal{P}_{S,F}} \mathcal{F}_t(i,j,k)\,dV
  = \int_{S}^{F} \left(\int_{\mathcal{P}_t} \mathcal{T}_i(t) \otimes \mathcal{T}_{ij}(t) \otimes \mathcal{T}_{ijk}(t)\,dV_t\right) dt
  = \mathcal{T}_i \otimes \mathcal{T}_{ij} \otimes \mathcal{T}_{ijk}
  = \langle \mathcal{P} \rangle.
\end{equation}
And consequently additional remarks and observations regarding tensor products and their representations over Riemannian Manifolds are discussed using the Fubini's Theorem and a proof by construction. Previous knowledge on Abstract Algebra, and Topology is assumed.\\ \\
\textbf{Keywords: } Topology, Abstract Algebra, Machine Learning.

%% ============================================================
\section{Introduction and Definitions:}
%% ============================================================

\lettrine[lines=1, lhang=0.1, loversize=0.2]{T}HE following definition of $\mathcal{P}$ is a p-dimensional Riemannian manifold product tensor with boundary. One may define such as:

\subsection{Definitions of a Tensor and a Tensor Set:}

Let a tensor $\mathcal{X}$ be defined as in:
\begin{equation}
  \mathcal{X} \subset \mathcal{T}_i \otimes \mathcal{T}_{ij} \otimes \mathcal{T}_{ijk}.
\end{equation}
For the following tensor $\mathcal{T}_i \otimes \mathcal{T}_{ij} \otimes \mathcal{T}_{ijk}.$ to be defined as in:
\begin{equation}
  \mathcal{T}_i \otimes \mathcal{T}_{ij} \otimes \mathcal{T}_{ijk}. \in \mathbb{T}^{+}.
\end{equation}
Where $\mathbb{T}^{+}$ is the set of all non-zero three-dimensional tensors such as $\mathcal{T}_i \otimes \mathcal{T}_{ij} \otimes \mathcal{T}_{ijk}.$

%% ============================================================
\subsection{Definitions of a Riemannian Manifold Product Tensor:}
%% ============================================================

\begin{definition}
\label{def:two-one}
    Let $\mathcal{P}$ denote a p-dimensional Riemannian manifold product with boundary. We will always assume $\mathcal{P}$ to factor three three-dimensional tensors that are non-zero. One may define such as:
    \begin{equation}
      \langle \mathcal{P} \rangle = \mathcal{X}.
    \end{equation}
    A general identity, defined as an integral over the subregion $\mathcal{P}_{S,F}$ with respect to the Riemannian volume form, where $t \in [S, F]$ is a foliation parameter indexing a family of hypersurfaces $\{\mathcal{P}_t\}_{t \in [S,F]}$ that foliates $\mathcal{P}_{S,F}$, and $i, j, k$ are continuous local coordinates on each leaf $\mathcal{P}_t$:
    \begin{equation}
      \langle \mathcal{P} \rangle_{g} = \int_{\mathcal{P}_{S,F}} \mathcal{F}_{t}(i, j, k) \, dV.
    \end{equation}
    Where $V$ is the Riemannian volume form on $\mathcal{P}$. And $g$ represents a discrete part of the variable $\mathcal{P}$. 
\end{definition}

%% ============================================================
\section{Additional Remarks regarding our Proofs:}
%% ============================================================

\lettrine[lines=1, lhang=0.1, loversize=0.2]{I}N this section, several remarks will be stated here regarding our observations of our theorems that have been written per Definition~\ref{def:two-one}.

\begin{remark}
\label{remark:one}
  There exist degenerate cases $\mathcal{P}_d$ of $\mathcal{P}$ for which the integral method does not converge in $\mathbb{R}$.
\end{remark}

%% ============================================================
\section{Proofs of Definition~\ref{def:two-one} over Riemannian Manifolds:}
%% ============================================================

%% ============================================================
\subsection{Definitions of Corollaries for Manifolds:}
%% ============================================================

\begin{corollary}
\label{cor:integrability}
  Let $\mathcal{F}_t$ be a continuous tensor field on the compact region $\mathcal{P}_{S,F} \subset \mathcal{P}$. Then $\mathcal{F}_t$ is integrable with respect to the Riemannian volume form $dV$, and $\int_{\mathcal{P}_{S,F}} \mathcal{F}_t(i,j,k)\,dV$ is well-defined.
\end{corollary}

\begin{theorem}
\label{theo:one}
  Per Definition~\ref{def:two-one}, the integral $\int_{\mathcal{P}_{S,F}} \mathcal{F}_t(i, j, k)\,dV$ defines a specific case of $\langle \mathcal{P} \rangle$ over a Riemannian manifold, where $\mathcal{F}_t$ is a continuous tensor field on $\mathcal{P}$ and $\mathcal{P}_{S,F}$ is the subregion bounded by regions $S$ and $F$. This follows from Corollary~\ref{cor:integrability} and Fubini's theorem for Riemannian submersions: if $\mathcal{P}_{S,F}$ admits a foliation by $t$-slices, the integral decomposes as an iterated integral over fibers, recovering the contraction $\mathcal{T}_i \otimes \mathcal{T}_{ij} \otimes \mathcal{T}_{ijk}$.
\end{theorem}

\begin{corollary}
\label{cor:tensor-nz}
  Per definitions in Definition~\ref{def:two-one}, since the tensors are assumed non-zero ($\langle \mathcal{P} \rangle$), there are non degenerate cases of Definition \ref{def:two-one} for both of the equations as restated below:
  \begin{equation}
    \langle \mathcal{P} \rangle = \mathcal{T}_i \otimes \mathcal{T}_{ij} \otimes \mathcal{T}_{ijk}.
  \end{equation}
  And thus:
  \begin{equation}
    \langle \mathcal{P} \rangle = \int_{\mathcal{P}_{S,F}} \mathcal{F}_{t}(i, j, k) \, dV.
  \end{equation}
\end{corollary}

\begin{corollary}
\label{cor:three}
  Per $\mathcal{F}_t(i, j, k)$ being factorizable as $\mathcal{T}_i(t) \otimes \mathcal{T}_{ij}(t) \otimes \mathcal{T}_{ijk}(t)$ on each leaf $\mathcal{P}_t$ of the foliation. There exist a conjugate variant of $\mathcal{P}_t$.
\end{corollary}

%% ============================================================
\subsection{Proof by construction regarding Corollaries and Definitions:}
%% ============================================================

\begin{proof}
\label{proof:construc}
We proceed by construction, invoking Corollaries \ref{cor:tensor-nz} and \ref{cor:three}.\\
(i) By Corollary~\ref{cor:tensor-nz}, the contraction $\mathcal{T}_i \otimes \mathcal{T}_{ij} \otimes \mathcal{T}_{ijk}$ is non-degenerate on $\mathcal{P}_{S,F}$, since the constituent tensors are assumed non-zero by Definition~\ref{def:two-one}. In particular, $\langle \mathcal{P} \rangle \neq 0$ on the initial leaf $\mathcal{P}_S$.\\
(ii) By Corollary~\ref{cor:three}, $\mathcal{F}_t(i,j,k)$ factors as $\mathcal{T}_i(t) \otimes \mathcal{T}_{ij}(t) \otimes \mathcal{T}_{ijk}(t)$ on each leaf $\mathcal{P}_t$. Suppose the identity $\langle \mathcal{P} \rangle = \int_{\mathcal{P}_{S,t}} \mathcal{F}_t(i,j,k)\,dV$ holds for all leaves up to $\mathcal{P}_t$. Since $\mathcal{F}_t$ is continuous and non-degenerate by Corollaries~\ref{cor:tensor-nz} and~\ref{cor:three}, the identity extends to $\mathcal{P}_{t+\epsilon}$ for any $\epsilon > 0$ by continuity of the foliation parameter.\\
(iii) Applying Fubini's theorem for Riemannian submersions over $\mathcal{P}_{S,F}$:
\begin{equation}
  \int_{\mathcal{P}_{S,F}} \mathcal{F}_t(i,j,k)\,dV
  = \int_{S}^{F} \left(\int_{\mathcal{P}_t} \mathcal{T}_i(t) \otimes \mathcal{T}_{ij}(t) \otimes \mathcal{T}_{ijk}(t)\,dV_t\right) dt
  = \mathcal{T}_i \otimes \mathcal{T}_{ij} \otimes \mathcal{T}_{ijk}
  = \langle \mathcal{P} \rangle.
\end{equation}
(iv) Since the factorization is non-degenerate on every leaf and Fubini assembles the leaves globally over $\mathcal{P}_{S,F}$, both equations of Definition~\ref{def:two-one} are simultaneously satisfied, completing the proof.
\end{proof}

%% ============================================================
\section{Conclusion:}
%% ============================================================

One may expand the observations and experimentation in Analysis and Complex Analysis. Moreover, per our abstract: We covered several products and their proofs over a Riemannian Manifold.

\nocite{*}
\bibliographystyle{acm}
\bibliography{pinter2010book, simons1968, weil1967basic, munkres}

%% ============================================================
\end{document}
%% ============================================================